\chapter{Transformers}
\label{ch:transformers}

\section{Mutual Inductance and Transformer Action}

A \textbf{transformer} consists of two (or more) coils --- a
\textbf{primary} and a \textbf{secondary} --- wound so that they share a
common magnetic flux, usually by winding both around the same core. A
changing current in the primary produces a changing flux (Faraday's law,
Chapter~\ref{ch:fields}); that changing flux links the secondary winding
and induces an EMF there. No direct electrical (galvanic) connection is
required between primary and secondary --- energy is transferred
entirely through the shared magnetic field, which is why transformers
are the standard way to provide \textbf{galvanic isolation}, in addition
to voltage or impedance transformation.

\section{The Ideal Transformer Equations}

For an ideal transformer (no losses, all flux shared between windings),
the ratio of secondary to primary voltage equals the ratio of secondary
to primary turns, the \textbf{turns ratio}:
\[
\boxed{\frac{U_2}{U_1} = \frac{N_2}{N_1} = n}
\]
Conservation of power ($U_1 I_1 = U_2 I_2$ for an ideal, lossless
transformer) then requires current to transform in the \emph{inverse}
ratio:
\[
\frac{I_2}{I_1} = \frac{N_1}{N_2} = \frac{1}{n}.
\]
A \textbf{step-up} transformer ($N_2 > N_1$) raises voltage and lowers
current; a \textbf{step-down} transformer ($N_2 < N_1$) does the
opposite. Neither creates or destroys energy --- power in equals power
out for an ideal transformer.

\begin{examplebox}[Turns ratio]
A transformer has 1000 turns on the primary and 50 turns on the
secondary, with 230\,V RMS applied to the primary. Find the secondary
voltage, and the secondary current if the load draws 2\,A on the primary
side.

\emph{Solution.} Turns ratio $n = 50/1000 = 0.05$, so
$U_2 = 230 \times 0.05 = 11.5\,\text{V}$. Current transforms inversely:
$I_2 = I_1/n = 2/0.05 = 40\,\text{A}$.
\end{examplebox}

\section{Impedance Transformation}

Because a transformer changes the voltage-to-current ratio, it also
changes the \emph{impedance} seen looking into the primary for a given
load impedance $Z_2$ on the secondary:
\[
\boxed{Z_1 = n^2 Z_2}
\]
This impedance-transforming property is exploited constantly in RF
work: audio output transformers match a low-impedance speaker to a
higher-impedance amplifier stage, and RF transformers (including
\textbf{ferrite-core broadband transformers} wound with twisted or
bifilar wire) match a transmitter's output impedance to an antenna
system, or match one stage's impedance to the next, without the
narrowband limitation of an LC matching network.

\begin{examplebox}[Impedance matching]
Find the turns ratio needed to match a 50\,$\Omega$ source to an
800\,$\Omega$ load.

\emph{Solution.} Require $Z_1 = n^2 Z_2 \Rightarrow n = \sqrt{Z_1/Z_2} =
\sqrt{50/800} = \sqrt{1/16} = 0.25 = 1{:}4$. A turns ratio of 1:4
(primary:secondary) steps the 800\,$\Omega$ load down to an apparent
50\,$\Omega$ at the primary.
\end{examplebox}

\section{Real Transformers: Losses and Efficiency}

Real transformers depart from the ideal model in several ways:

\begin{itemize}
\item \textbf{Copper losses}: $I^2R$ heating in the winding resistance,
worse at high current (heavy secondary loading).
\item \textbf{Core losses}: hysteresis loss (energy needed to repeatedly
reverse the core material's magnetic domains) and eddy-current loss
(circulating currents induced within the core itself), both of which
increase with frequency --- one reason mains-frequency (50/60\,Hz) power
transformer cores cannot simply be reused, unmodified, at radio
frequencies.
\item \textbf{Leakage inductance}: flux produced by one winding that
does \emph{not} link the other winding, behaving as an unwanted series
inductance.
\item \textbf{Winding capacitance}: parasitic capacitance between turns
and between windings, which becomes significant at high frequency.
\end{itemize}

Efficiency is defined the same way as for any energy-conversion device:
\[
\eta = \frac{P_{\text{out}}}{P_{\text{in}}} \times 100\%.
\]
Well-designed power transformers commonly exceed 95\% efficiency at
rated load; efficiency generally drops at light load (core losses become
a larger fraction of a small output) and at heavy overload (copper
losses dominate and the core may begin to saturate).

\section{Autotransformers}

An \textbf{autotransformer} uses a single tapped winding rather than
electrically separate primary and secondary windings. It is smaller,
cheaper, and more efficient than an equivalent two-winding transformer
for a given power rating, but sacrifices galvanic isolation between
input and output --- an important safety distinction from a conventional
isolation transformer. Variable autotransformers (e.g., the classic
``Variac''-type device) are commonly used on the test bench to control AC
line voltage.

\section{RF Transformers and Baluns}

At radio frequencies, transformers are frequently wound on ferrite
toroid or binocular cores using twisted-pair or bifilar windings to
achieve tight coupling (high efficiency) over a wide bandwidth. A
particularly important RF application is the \textbf{balun}
(``balanced-to-unbalanced'' transformer), used to connect an unbalanced
coaxial feedline to a balanced antenna element such as a dipole
(Chapter~\ref{ch:antennas}) without upsetting the antenna's current
balance --- a common source of feedline radiation and RF-in-the-shack
problems if omitted or poorly chosen. A related device, the \textbf{unun}
(``unbalanced-to-unbalanced''), transforms impedance between two
unbalanced points, commonly used to match a random-wire or end-fed
antenna to a 50\,$\Omega$ coaxial feedline.

\begin{quickreview}
\item Transformer action relies on mutual inductance: a changing primary
current induces an EMF in a magnetically coupled secondary.
\item $U_2/U_1 = N_2/N_1 = n$; current transforms inversely, $I_2/I_1 =
1/n$; ideal power in equals power out.
\item Impedance transforms as $Z_1 = n^2 Z_2$ --- the basis of
transformer impedance matching.
\item Real transformers have copper loss, core loss (hysteresis and eddy
current), leakage inductance, and winding capacitance, all limiting
efficiency and bandwidth.
\item RF baluns and ununs are specialized transformers for connecting
balanced/unbalanced feedlines and antennas without disturbing current
balance.
\end{quickreview}

\begin{practiceproblems}
\item A transformer steps 120\,V down to 12\,V. If the primary turns
count is 600, find the secondary turns count.
\item Find the primary current if the secondary of the above transformer
delivers 3\,A to a load (assume ideal, lossless transformer).
\item Find the turns ratio needed to match a 50\,$\Omega$ transmitter
output to a 200\,$\Omega$ antenna feedpoint impedance.
\item A transformer delivers 90\,W output while consuming 100\,W input.
Find its efficiency, and name two loss mechanisms that could account for
the difference.
\item Explain why a balun is typically used where a coaxial feedline
connects to a dipole antenna, in terms of balanced vs.\ unbalanced
currents.
\end{practiceproblems}
